\paragraph{Limitations and next steps.} The inclusion of behavior in representation is a near-definitional observation for a probe that contains the system's procedure, whereas the evidence that the inclusion is strict comes from the systems we tested. Our tests also leave attribute binding, the modifier--head relation that is the main category of T2I-CompBench \citep{huang2023t2i,rassin2022}, as an untested relation type, and the random-hierarchy result of Appendix~\ref{app:extra} tests only the input side, showing that the one-hot input of that grammar recovers the latents of new symbols poorly, while trained networks on that grammar remain untested. Generative vision--language models, in-context learning and chain-of-thought prompting also remain untested, and the paired placement predicts that they inherit positional binding wherever their states follow word order or location. Our evidence for design comes from repairs and from a recovery objective that also adds data, so separating the objective from the data is the next test. Our trained networks are small, with rendered domains so that we could keep parts new to them, and their recoverability comes largely from their inputs. In pretrained models parts are new only to the probe, apart from invented words. Additions without forgetting, a new part in every role, nesting and code swaps rest on toy systems alone, and testing them in pretrained models, for instance by writing an invented word's vector into a role or by nesting relations in a sentence, is the next step. Fixed quadrants tie each role to a location, and our relative placement loosens that tie without removing it, so a canvas on which position carries no information about the role would settle how much of the role the encoder states mark. Image encoders still lack a tested repair. The harness repairs of Appendix~\ref{app:harness} rest on lexicons of the training relations, on a planner that we prompt with four demonstrations and on GLIGEN's boxes, and a harness whose decomposer must learn held-out relation words remains untested. FLUX, a gated model, also remains untested among the transformer denoisers. Our probe of the denoiser on its own generated images, fitted on 344 boxes, names the nouns only weakly and cannot tell figure from ground, so whether the denoiser keeps the places it fails to draw remains open. Richer methods, such as a discovered tensor-product decomposition \citep{mccoy2019}, may find structure that our probes miss. A human study could ask whether people show the same gap, by testing whether they recover the parts of wholes built from newly learned parts.
